% Document class `report-template` accepts either project-plan or final-report option in
% []. This will change the title page as necessary.
% \documentclass[project-plan]{report-template}
\documentclass[final-report]{report-template}
\usepackage{multirow}


% Packages I use in my report.
\usepackage{graphicx}
\usepackage{amsmath}
\usepackage{blindtext}
\usepackage{array}  
\usepackage{caption}
\usepackage{float}
\usepackage{rotating}


% Directory where I saved my figures.
\graphicspath{{./figures/}}

% Metadata used for the title page - please modify.
\university{Imperial College London}
\department{Department of Earth Science and Engineering}
\course{MSc in Applied Computational Science and Engineering}
\title{OpenWeather ChatAgent: Design and Implementation of an LLM-Based Weather Assistant System}
\author{Fangze Liu}
\email{fangze.liu24@imperial.ac.uk}
\githubusername{esemsc-fl724}
\supervisors{Prof. Christopher Pain\\
             Dr. Olga Buskin}
\repository{https://github.com/ese-ada-lovelace-2024/irp-fl724}

\begin{document}

\maketitlepage  % generate title page

% Table of Contents page
\tableofcontents
\newpage   % Break page, separate page after table of contents

\section*{AI Acknowledgement Statement}
I used OpenAI’s ChatGPT‑4o (https://chat.openai.com/) to support my learning and understanding of large language models (LLMS), as well as to debug and test the code of my projects. This generative AI tool supported my learning process, but the submitted work is my own and it reflects my own understanding and effort as I declared by signing the Academic Integrity Declaration.

% Abstract
\section*{Abstract}
This project aims to design and implement an AI weather agent system powered by Large Language Models (LLMs). By integrating with the OpenWeather API, the system supports natural language queries for real-time weather. With the help of OpenAI's Function Calling, the model can dynamically interpret the user's intention and structurally call external tool functions to obtain weather data. To enhance modularity and cross-model compatibility, the project will explore the use of the Model Context Protocol (MCP). Optional extensions may include multi-model support and lightweight memory modules. The agent will be optimized in terms of intelligence, stability, and responsiveness. The final deliverables of the project include a user interface, end-to-end system integration, and full documentation. This project aims to demonstrate how LLMs can evolve from passive language generators into AI agents, providing a practical showcase of LLM-powered services in real-world applications.

% Introduction section
\newpage
\section{Introduction}
\subsection{Problem Description and Significance}
Weather, as the most critical environmental information after time and place, has a direct and far-reaching impact on people's daily lives. However, traditional APIs or weather query systems mostly use template-based responses or keyword matching, which are unable to understand the diversified expressions of user's natural language, and are difficult to adapt to the increasingly complex query needs of users. 

In recent years, with the rapid development of LLMs in natural language processing, more and more users have started to use conversational LLMs like ChatGPT (OpenAI, 2023)~\cite{OpenAI2023} for various queries. However, many user questions require knowledge beyond what a pre-trained LLM can store internally.~\cite{Zhang2023} For example, a user might ask, "What’s the weather like in London right now?" Clearly, such tasks require accessing real-time data through external APIs.~\cite{Chadha2024} The ability to interact with tools is a key step in the evolution of LLMs from passive text generators to action-capable AI agents.

Nevertheless, applying LLMs to real-world weather queries presents several challenges, such as: how to understand the user's intention, how to call the APIs in a structured way to obtain the external data; and how to semantically process the returned data.

Building an AI Weather Assistant System based on Large Language Models (LLMs) not only can enhance the user experience, but also create a practical example for engineering applications for LLMs. These design patterns can be applied to other agent-based cases such as tourism, health and education etc.


\subsection{Review of Existing Work}
In my research, I have discovered several different methods for designing weather assistant systems.

First, there are traditional weather information systems, which usually rely on APIs and mobile applications to provide weather data~\cite{Zhang2023}. Users must manually interpret numerical weather data, which is not user-friendly, especially in varied real-life situations.

Second, some weather chatbots are built using the Rasa framework~\cite{WeatherBot2018}, which use a combination of traditional NLP, predefined rules, and API calls to deliver responses. However, the disadvantage is that it relies on predefined intents and rules, making it difficult to handle complex real-world questions.

In the past two years, AI agents have become more popular. Developers can use system prompts to tell LLMs what tools are available in the system. Some weather systems~\cite{Jeong2024} have explored the integration of LLMs with weather APIs. 

However, there are still some limitations. For example, the tool lists are directly embedded in the LLM prompt, making it hard to scale and maintain. Also, due to the lack of a clear standard, LLM may return incorrect formats, causing retries and extra token usage. Some well-known agents, such as Cline~\cite{Cline2025}, are using the system prompt method.

After recognizing the problems with system prompts, major model providers started to introduce a more structured way of calling tools, known as Function Calling~\cite{Liu2024}~\cite{Gim2024}. Platforms like OpenAI~\cite{OpenAI2025}, Claude~\cite{Claude2025}, and Gemini~\cite{Gemini2025} have all developed their own Function Calling systems. The core idea is to register each tool using a unified JSON Schema format, and provide that to the model separately from the prompt. I have also seen some articles~\cite{Anthonio2024}~\cite{KermaniWeatherAgent2025} on implementing weather chatbot that have already started using Function Calling.

Function Calling has its benefits: developers do not need to write complex prompts, the system can detect format errors and recover, and it also saves tokens by not repeating tool descriptions. However, the main drawback is the lack of a unified standard. Different providers define Function Calling differently, and some models do not natively support this feature, making tool definitions difficult to reuse.

As a result, both System Prompt–based and Function Calling–based methods for defining tools currently coexist in the market.

To address these issues, providers begin to explore some methods for managing tools, such as the Model Context Protocol~\cite{Anthropic2024MCP}, which will be discussed in detail in later sections.

\subsection{Objectives}
This project aims to develop a practical intelligent weather assistant system based on large language models (LLMs), supporting natural language interaction, real-time weather data queries, and personalized weather reports.

The system is designed to integrate multiple modules, including intelligent location, weather query, stylized output and memory modules, etc. The system supports multilingual conversations, multiple output styles, and user preference recognition. It also supports short-term and long-term memory. To enhance multi-model compatibility and maintainability, this project adopts the Model Context Protocol (MCP) for unified tool registration and combines Retrieval-Augmented Generation (RAG)~\cite{Jeong2024}~\cite{GaoYunfan2023} for knowledge enhancement.

The following deliverables are expected from the project:

\begin{itemize}
  \item A working intelligent weather chatbot system, supporting both command-line and web-based interaction, integrated with LLMs and the OpenWeather API~\cite{OpenWeatherAPI}.
  \item A modular toolset based on the MCP protocol, supporting services such as weather, location and map visualization.
  % \item A series of personalized features for system expansion.
  \item Evaluate the feasibility of the method and quantify the improvement in performance.
  % \item A complete project report, source code, demonstration video, and deployment documentation for reproducibility and delivery.
\end{itemize}



% Another section
\newpage
\section{Methodology}
\subsection{System Architecture and Technology Stack}

The system for this project adopts a microservice architecture and uses the Model Context Protocol (MCP) to manage various tools and services. The core framework is built on a combination of LLM, MCP, and RAG, which shows the popular trends in AI Agent. The system is divided into four core layers, each with clear responsibilities and interface definitions, in accordance with back-end development standards.


\textbf{LLM-driven layer (Agent):}
This layer uses models such as OpenAI GPT~\cite{OpenAI2025} or DeepSeek~\cite{DeepSeek2025} as the main reasoning engine. The main reason for this choice is that it supports Function Calling, which simplifies tool invocation. This layer is responsible for understanding the intentions of the users, planning the execution steps, and coordinating calls to external tools. The entire reasoning process is implemented using LangGraph~\cite{LangGraph2025}, following ReAct-style process, which has good traceability and execution control capabilities.

\textbf{MCP tools layer (MCP server):}
It is the core integration layer of the system. Weather-related tools are mainly managed through the MCP protocol. This layer supports dynamic discovery, registration, and execution of tools, providing a standardized tool interface for LLM agents.

\textbf{RAG knowledge layer:}
Using a locally deployed ChromaDB vector database~\cite{Chroma2025}, build collections such as weather knowledge and city information. Combine the Sentence Transformers embedding model to achieve fast semantic retrieval and knowledge enhancement, thus improving the accuracy of system responses.


\textbf{User interface layer:}
This project separates the front-end and back-end. The front-end is responsible for calling the back-end interface. The front-end uses Streamlit~\cite{Streamlit2025} to implement a modern web interface that supports streaming output. The interface is visually appealing and does not require users to operate on the command-line, improving the user experience.

\subsection{Model Context Protocol (MCP)}
\subsubsection{The Dilemma of Function Calling}
As mentioned in the introduction of this paper, Function Calling enables large language models to use external tools by creating an intermediary external function that passes requests from the LLM and executes the corresponding tool call.

However, this technology also has a significant drawback: the code implementation workload is substantial. Even a simple external function can require hundreds of lines of code. To make the LLM understand these functions, developers must manually write a function specification in JSON Schema format for each external tool and design a prompt template description. This development model forces developers to build custom adaptation interfaces for each system, making the development process time-consuming, labor-intensive, and difficult to maintain in the long term. 

\subsubsection{Model Context Protocol: Architecture and Advantages}
To address the “Tower of Babel dilemma” caused by Function Calling, Anthropic introduced the Model Context Protocol (MCP) in November 2024. MCP is not a specific tool but a standardized protocol for AI agent development. It adopts a distributed architecture, separating the client (LLM-driven) from the server (external tool provider). With this structure, a single client implementation can access any MCP-compliant server, significantly improving development efficiency.


The main advantages of MCP include:

\begin{itemize}
    \item \textbf{Standardization:} MCP defines a unified protocol for tool invocation, including interface definitions, parameter formats, error handling, and documentation specifications. As a result, MCP tools can be directly reused across different projects without additional adaptation.

    \item \textbf{Modular Architecture:} Unlike the tightly coupled design of Function Calling, MCP separates tool providers from tool consumers and supports asynchronous execution, offering better scalability and parallelism.

    \item \textbf{Ecosystem:} MCP offers SDKs for Python, TypeScript, and Java, allowing developers to build clients and servers with just a few lines of code. This lowers the barrier to entry and encourages rapid development.
\end{itemize}

\subsubsection{Development Difficulty Comparison}

This project initially adopted a Function Calling-based Agent system and was later upgraded to an MCP-based Agent baseline. 

To evaluate the advantages of the MCP architecture in terms of development efficiency, the project compared the code complexity of Function Calling and MCP under the same functionality. When supporting 10 tools, Function Calling requires approximately 150 lines of code, while MCP requires only 50 lines, representing a 67\% reduction in code volume. Additionally, when adding a new tool, Function Calling requires approximately 15 lines of schema configuration, whereas MCP only needs about 4 lines of decorators, significantly reducing the cost of scalability.

Therefore, this project will use the MCP protocol unified management tool to support automatic discovery and registration of modules such as location parsing, weather queries, air quality, and output formatting.

\subsection{Agent Design}
AI agent is a software system that can perceive its environment, make decisions, and take actions to achieve specific goals. It uses artificial intelligence techniques to act on behalf of users to pursue objectives and complete tasks. In this project, LLM serves as the core of the agent design.

\subsubsection{Reasoning Strategy}
This project follows the ReAct (Reasoning and Acting) process, forming a closed-loop process of reasoning, acting, and observation:

\begin{itemize}
    \item \textbf{Reasoning phase:} The LLM analyzes user input and develops a tool invocation plan based on the context.
    \item \textbf{Execution phase:} The system sequentially calls MCP tools, processes the results, and handles potential errors.
    \item \textbf{Observation phase:} Intermediate states are collected to support subsequent reasoning adjustments and memory updates.
    \item \textbf{State Management:} Contextual memory is maintained using the MemorySaver module from the LangGraph framework.  
\end{itemize}

\subsubsection{Single-Agent vs Multi-Agent}
Currently, Agent architectures can be broadly categorized into single-agent and multi-agent structures.

A Single-Agent System integrates perception, planning, and tool execution into a single loop driven by an LLM. The workflow is straightforward: the agent receives a task, invokes the required tools via MCP, and produces an integrated response based on the outputs.

A Multi-Agent System (MAS) is a system composed of multiple autonomous agents that can communicate, coordinate, and collaborate to complete complex tasks. MAS adopts an “expert team” model, where each agent specializes in specific tasks or domains. There are various MAS organizational structures, such as the Supervisor–Expert architecture, Hierarchical structure, and Blackboard system. 

If a single-agent system with MCP resembles an “all-round expert” capable of using various tools independently, then a multi-agent system is more like a “team of experts”, where each member handles the domain they excel at.

\subsubsection{Experimental Evaluation of Agent Architectures
}
This project conducted experimental evaluations of agent architecture selection. The system used multi-agent systems based on single-agent architecture and supervisor-expert architecture to build weather assistants, and compared and analyzed their performance. The multi-agent system consists of four experts, each responsible for location, weather, air quality, and reporting.

\begin{table}[ht]
\centering

\begin{tabular}{|c|c|c|c|}
\hline
\textbf{Model} & \textbf{Metric} & \textbf{Single Agent} & \textbf{Multi-Agent} \\
\hline
\multirow{4}{*}{GPT-4o-mini} 
& Total Response Time (s)     & 4.62   & 12.14 \\
& Token Usage           & 5646   & 11875 \\
& API Calls             & 2.0    & 9.33 \\
& Execution Success     & 100\%  & 100\% \\
\hline
\multirow{4}{*}{DeepSeek-V3} 
& Total Response Time (s)     & 12.12  & 33.45 \\
& Token Usage           & 9275   & 21059 \\
& API Calls             & 2.0    & 9.67 \\
& Execution Success     & 100\%  & 98\% \\
\hline
\end{tabular}
{\small\caption{Cross-Model Comparison of Single-Agent and Multi-Agent Performance: Averaged Results Over 10 Runs on 5 Weather-Related Queries}}
\end{table}

The experimental results show that in the context of this project, compared with single agents, multi-agent architectures have longer response times and higher token consumption in single-round queries, and do not show any obvious advantages in terms of success rate.

Based on the above design and experiments, this project determined that the single-agent architecture centered on LLM and MCP should continue to be improved. This architecture enhances the flexibility and robustness of dialogue, laying the foundation for personalized responses and multi-round interactions.

\subsection{Intelligent Module Implementation}
This system can be divided into multiple modules based on functionality, each designed for different tasks, such as user intent analysis, user profiling, weather data acquisition, and personalized output. These modules are composed of different tools and are uniformly invoked through the MCP protocol to collaborate and run in the agent. The modules follow a clear data flow sequence and optimize performance and response efficiency through caching and asynchronous execution mechanisms.

\subsubsection{User Profiling and Memory System}
The goal of this system is to create a personalized weather assistant, so implementing a memory mechanism is crucial. For an intelligent chatbot, short-term memory is a basic requirement. The system must remember certain contextual information from the current conversation to support users' subsequent operations based on previous questions. On the other hand, long-term memory is essential for achieving personalization. Different users have significant differences in their needs for weather services, such as preferred weather report styles, temperature units used, frequently queried cities, and even language and emotional expression habits.

In summary, this system is primarily divided into two layers of memory, supporting a hierarchical structure for short-term and long-term memory. Short-term memory records the context and temporary preferences of the current conversation, while long-term memory persistently stores user interests, preferences, and other data in JSON format, supporting cross-conversation memory. The system intelligently merges short-term and long-term memory, persistently saving these user characteristics. This enables the system to continue and adapt to users' historical preferences in subsequent interactions, significantly enhancing the user experience and service consistency.

The system's extraction of user preferences is divided into explicit identification and implicit identification. User preferences are extracted using regular expressions, while implicit identification occurs when a user maintains a habit beyond a certain threshold, at which point the system automatically upgrades it to a preference and saves it as part of the user profile.

\subsubsection{Location Resolution Module}
The primary task of this module is to convert the location to be queried into coordinates for use with the weather API. The system supports multiple input mode recognition, including city names, postal codes, GPS coordinates, ZIP codes, and country codes, using regular expressions to automatically determine the input type.

This module employs a dual API fallback strategy. It primarily calls the OpenWeather API's Geocoding interface, which supports geocoding, ZIP code parsing, reverse geocoding, and more. If the call fails or the returned information is incomplete, it transparently falls back to the more specialized location query API, OpenCage, and merges the supplementary fields into the final result, enabling the system to provide more comprehensive geographic information data.  

Additionally, the system incorporates a caching mechanism based on hash cache keys, effectively avoiding duplicate parsing requests and improving response speed and API call efficiency.  

\subsubsection{Weather-Related Services}

To ensure high-quality responses to user queries, the system implements a comprehensive set of weather-related service tools, covering weather data retrieval, air quality monitoring, alert notification, and weather map visualization. All data is sourced from the OpenWeather API, ensuring accuracy and reliability.

These modules together form a high-throughput pipeline: \\
Location Parsing → Concurrent Weather Data Retrieval → Profile Fusion → Stylized Output Generation.

\paragraph{Weather Data Module}~

This module is built on OpenWeather APIs (versions 2.5 and 3.0) and serves as the foundation for acquiring essential meteorological information. It supports current weather, minutely forecasts, hourly forecasts, 16-day daily forecasts, and historical weather data. Users can query by city name, geographic coordinates, or postal code. \\
The module also offers basic climate statistics, such as average temperature, humidity, wind speed, and precipitation frequency over a specified period, which helps users analyze weather trends. \\

A hierarchical caching mechanism is designed for this module, assigning differentiated cache durations to various data types to balance freshness and reduce API load. To enhance stability, it incorporates an exponential backoff retry mechanism and supports Gzip-compressed data transmission.

\paragraph{Air Quality Module}~

Inheriting the structure of the weather module, this component uses the OpenWeather Air Pollution API to retrieve pollutants such as PM2.5, PM10, O$_3$, NO$_2$, SO$_2$, and CO. It also calculates the unified Air Quality Index (AQI) to assess the level of air pollution and its potential health risks.

\paragraph{Weather Alert Module}~

This module integrates both OpenWeather’s official disaster alert API and user-defined alert conditions, forming a multi-layered risk notification system. Official alerts are automatically categorized and prioritized, with multilingual output supported. \\
In parallel, users can define custom thresholds for critical indicators (e.g., temperature, wind speed, AQI). Once these thresholds are exceeded, the system triggers personalized alerts and provides tailored recommendations based on the user profile (e.g., for children, the elderly, or drivers).

\paragraph{Weather Map Module}~

This module utilizes the OpenWeather Maps API to render a variety of meteorological layers (temperature, wind speed, cloud cover, precipitation, etc.) as well as specialized layers (e.g., PAC0, TA2, HRD0). It supports customization of layer transparency, terrain overlays, and solar illumination. \\
The system automatically stitches tile images across geographic areas to generate weather visualizations for specified regions, offering users an intuitive graphical overview of atmospheric conditions.

\subsubsection{Personalized Weather Reporting Module}

To deliver personalized weather reports, the system supports stylized language wrapping of raw meteorological data. A keyword-based style recognition mechanism enables over 20 output styles, including brief, detailed, social media, poetic, and professional tones (e.g., pilot, scientist), among others.

There are two modes of output generation:

\begin{itemize}
  \item \textbf{Template-Based Output:} \\
  Common styles are handled using predefined system prompts. These allow for direct and fast rendering of weather information.

  \item \textbf{LLM-Driven Output:} \\
  For complex styles that require greater expressiveness, the system injects structured weather data and user profile information into a customized prompt, guiding the LLM to generate content in the target style. \\
\end{itemize}

All formatted outputs are cached using an MD5-based hashing mechanism to avoid redundant LLM calls and reduce resource consumption.

\subsection{Retrieval-Augmented Generation (RAG) for Knowledge Enrichment}

To strengthen the system’s domain accuracy and contextual reasoning, we incorporated a Retrieval-Augmented Generation module. By building a structured weather knowledge base, refining the embedding and search pipeline, and injecting retrieved facts on-the-fly during dialogue, the model delivers more complete and reliable answers to complex weather queries.

\subsubsection{Knowledge-base Construction}

\textbf{Data collection:} Twelve thematic corpora were assembled from Wikipedia and curated documents, covering meteorology fundamentals, health impacts, hazard mitigation, historical weather records, and global city profiles (climate, culture, transport, cuisine, etc.). Raw material is stored as JSON ($\approx$ 4.3 MB, 73,373 rows).

\vspace{0.3em}
\textbf{Vector store:} Each document is split with a recursive character splitter to minimise semantic fragmentation, then embedded with Sentence-Transformers \texttt{all-MiniLM-L6-v2} (384-D, FP16). Embeddings and metadata (title, source, topic) are written to a local ChromaDB instance, enabling topic-aware re-ranking. Keyword filters and relevance scoring eliminate off-topic or low-quality passages before ingestion.

\subsubsection{Embedding Strategy and Chunk-size Optimisation}

Chunk size was tuned experimentally. With 300 tokens per chunk, the system achieves the best balance of accuracy ($\approx$ 85\%), knowledge coverage ($>$ 93\%) and latency ($\approx$ 1.2 s). Smaller chunks ($\leq$ 200 tokens) lower latency but harm recall, while larger chunks ($\geq$ 500 tokens) slow retrieval without significant accuracy gains. Accordingly, the default setting is 300 tokens with zero overlap, embedded via the MiniLM model described above.

\begin{figure}[H]
    \centering
    \includegraphics[width=1.0\textwidth]{latex/figures/all_similarity_comparison.png}
    \caption{Comparison of Mean Similarity Across Different Chunking Strategies}
    \label{fig:all_similarity_comparison}
\end{figure}

\subsubsection{Retrieval and Intelligent Routing}

Search employs a hybrid scheme. A cosine top-k query ($k \leq 5$) brings back candidate chunks from ChromaDB; keyword overlap and other features are then linearly combined into a final score. An intent classifier routes each user question to the most relevant sub-collection (general weather, city facts, hazard advice, and so on), further improving match quality.

\subsubsection{Knowledge Fusion and Contextual Enhancement}

The RAG layer works in parallel with the user-profile component. Before formulating a reply, the system builds a prompt that blends:
\begin{itemize}
    \item conversational context,
    \item user preferences --- style, language, focus city,
    \item and retrieved knowledge snippets.
\end{itemize}
The same mechanism auto-fills missing information in elliptical queries, for example by inserting the current city or metric, so responses remain coherent and personalised.

\subsubsection{Module Benefits and Evaluation}

Benchmarking shows that, on multilingual, cross-domain questions rich in technical terms, the RAG-enabled assistant consistently outperforms a plain LLM baseline. Gains are observed in factual accuracy, semantic depth, and completeness of explanation.


\subsection{Performance Optimization and Engineering Trade-offs}









\section{Results}

\section{Discussion}

\section{Conclusions}



% References
\newpage
\bibliographystyle{plain}
\bibliography{references}  % BibTeX references are saved in references.bib

\end{document}          

